\subsection{切片分解。局部同胚下测度的逆像}

设 X 是一个局部紧空间，\(\pi\) 是一个由 X 到局部紧空间 T 中的映射，\(\mu\) 是 T 上的一个正测度，\(\Lambda: t \mapsto \lambda_t\) 是一个标量本质 \(\mu\) -可积且淡 \(\mu\) -可测的映射（由 T 到 \(\mathcal{M}_{+}(X)\) 中）。设 \(\nu = \int \lambda_t d\mu(t)\) 。若 \(\lambda_t\) 由 \(\overline{\pi}(t)\) 承载（对每个 \(t \in T\) ），则称等式 \(\nu = \int \lambda_t d\mu(t)\) 为 \(\nu\) 关于 \(\pi\) 的一个切片分解（或一个分解）。这一概念将在 Ch. VI 中详细研究。

命题 10. --- 沿用上述记号，设 \(\pi\) 是 \(\nu\) -可测的。设 \(g\) 是 T 上的函数 \(t \mapsto \lambda_t^*(1)\) 。为使 \(\pi\) 是 \(\nu\) -逆紧的，必须且只须 \(g\) 是局部 \(\mu\) -可积的，此时

\[
\pi (\nu) = g \cdot \mu .\tag{14}
\]

我们先在 g 局部 \(\mu\) -几乎处处有限的假设下进行论证；将在证明末尾去掉这一辅助假设。由于 \(\pi\) 依假设是 \(\nu\) -可测的，说 \(\pi\) 是 \(\nu\) -逆紧的等价于说 \(\nu^{\bullet}(f\circ\pi)<+\infty\) 对每个函数 \(f\in\mathcal{K}_{+}(\mathrm{T})\) 成立；g 既然局部几乎处处有限，我们就处于可应用 §3, No.~2 的命题 5 的断言 c) 的条件之下。因此

\[
\int^ {\bullet} (f \circ \pi) d \nu = \int^ {\bullet} d \mu (t) \int^ {\bullet} (f \circ \pi) d \lambda_ {t} = \int^ {\bullet} f (t) g (t) d \mu (t),
\]

从 \(\lambda_{t}\) 集中于 \(\overline{\pi}^{1}(t)\) 这一事实得出。我们知道 g 是 \(\mu\) -可测的，因为 \(\Lambda\) 是 \(\mu\) -适宜的（\(\S3\) , No.~1, 定义 1）。于是，说第一个成员对每个 \(f \in \mathcal{K}_{+}(\mathrm{T})\) 为有限，就等价于说 g 是局部 \(\mu\) -可积的（\(\S5\) , 命题 1），这时 (14) 由上述关系式立即得出。

于是只剩下消去辅助假设。若 g 是局部 \(\mu\) -可积的，则 g 局部 \(\mu\) -几乎处处有限，该假设确实满足。设 \(\pi\) 是 \(\nu\) -逆紧的，我们来证明 g 局部几乎处处有限。设 \(\mathfrak{K}\) 是 \(\mu\) -稠密的紧集 K（使 \(\Lambda|K\) 为淡连续者）所成之集；由于 g 可测，我们归结为证明每个紧集 \(K \in \mathfrak{K}\) （满足 \(g|K = +\infty\) ）都是 \(\mu\) -可忽略的。现设 \(\mathcal{H}\) 是函数 \(h \in \mathcal{K}_{+}(X)\) （满足 \(h \leqslant 1\) ）所成的集；置 \(g_{h}(t) = \lambda_{t}(h)\) ，用 \(\Lambda_{h}\) 表示 \(\mu\) -适宜映射 \(t \mapsto h \cdot \lambda_{t}\) ，用 \(\nu_{h}\) 表示 \(\Lambda_{h}\) 的积分，并用 f 表示 \(\mathcal{K}_{+}(T)\) 中满足 \(f \geqslant \varphi_{K}\) 的一个元素。对确实满足辅助假设的 \(\Lambda_{h}\) 应用公式 (14)，得到：

\[
\int (f \circ \pi) d \nu \geqslant \int (f \circ \pi) d \nu_ {h} = \int f g _ {h} d \mu .
\]

但函数 \(fg_h|K\) 在 \(K\) 上构成一个递增有向的连续函数集，其上包络取值 \(+\infty\) ；由 Dini 定理（GT, X, §4, No.~1, 定理 1），可以选取 \(h\) ，使得 \(fg_h|K\) 大于或等于任意一个正数 \(n\) ，由此得出 \(\int(f \circ \pi) d\nu \geqslant n \mu(K)\) 。由于 \(\pi\) 是 \(\nu\) -逆紧的，第一个成员有限，于是推出 \(\mu(K) = 0\) 。

推论 1. --- 设 \(\pi\) 是 \(\nu\) -可测的。

\begin{enumerate}
\def\labelenumi{\alph{enumi})}
\item
  若 \(\mathrm{N}\subset \mathrm{T}\) 是局部 \(\mu\) -可忽略的，则 \(\overline{\pi}^{1}(\mathrm{N})\) 是局部 \(\nu\) -可忽略的。
\item
  若 \(f\) 是一个 \(\mu\) -可测映射（由 \(T\) 到拓扑空间 \(G\) 中），则 \(f \circ \pi\) 是 \(\nu\) -可测的。
\end{enumerate}

我们重新采用前述证明末尾的记号 \(\Lambda_h\) , \(\nu_h\) , \(g_h\) ：由于 \(\nu_h\) 对每个 \(h \in \mathcal{H}\) 是有界测度，故 \(\pi\) 是 \(\nu_h\) -逆紧的，\(g_h\) 是局部 \(\mu\) -可积的，且 \(\pi(\nu_h) = g_h \cdot \mu\) 是一个以 \(\mu\) 为基的测度。由此推出，N （相应地，\(f\) ）对测度 \(\pi(\nu_h)\) 是局部可忽略的（相应地，可测的）（\(\S 5\) , No.~3, 命题 3 的推论 1 与命题 4）。因而 \(\overline{\pi}^1(N)\) （相应地，\(f \circ \pi\) ）对测度 \(\nu_h\) 是局部可忽略的（相应地，可测的）（命题 2 的推论 2，相应地为命题

3）。最后，注意到诸测度 \(\nu_{h}\) 构成一个递增有向的正测度族，其上确界为 \(\nu\) ，并应用 §1, No.~4 的命题 11 的推论 1（相应地，推论 2）。

推论 2. --- 设 \(\pi\) 是 \(\nu\) -逆紧的；设 \(\mathbf{f}\) 是定义在 \(\mathbf{T}\) 上、取值于一个 Banach 空间或 \(\overline{\mathbf{R}}\) 中的映射。为使 \(f \circ \pi\) 是本质 \(\nu\) -可积的，必须且只须 \(gf\) 是本质 \(\mu\) -可积的。

注意到命题 10，这由 §5, No.~3 的定理 1 以及 No.~2 的定理 1 立即得出。

例. --- 设 X 和 T 是两个局部紧空间，并设 \(\pi\) 是 X 到 T 中的一个局部同胚。换句话说（GT, I, §11, 习题 25），我们假设每个点 \(x \in X\) 有一个邻域 V，使得 \(\pi|V\) 是 V 到 \(\pi(x)\) 的一个邻域上的同胚；必要时把 V 换成 x 的一个相对紧开邻域 W，使得 \(\overline{W} \subset V\) ，便可推出：集 \(\mathcal{U}\) （由 X 的相对紧开子集 U 中使得 \(\pi|\overline{U}\) 是 \(\overline{U}\) 到其像上的同胚者组成）是 X 的一个开覆盖。现设 \(\mu\) 是 T 上的一个正测度；若 U 是 \(\mathcal{U}\) 的一个元素，则 \(\pi(U)\) 是紧空间 \(\pi(\overline{U})\) 中的一个开集，因而是 T 的一个局部紧子空间，并且我们知道如何定义测度 \(\mu|\pi(U)\) （由 \(\mu\) 在 \(\pi(U)\) 上诱导的）（Ch. IV, §5, No.~7）。设 \(\nu_{U}\) 是 \(\mu|\pi(U)\) 在 \(\pi|U\) 的逆同胚下的像；我们将证明存在一个且仅一个 X 上的测度 \(\nu\) ，它诱导出测度 \(\nu_{U}\) （在每个开集 \(U \in \mathcal{U}\) 上）。这个测度称为 \(\mu\) 在局部同胚 \(\pi\) 下的逆像，记作 \(\overline{\pi}^{1}(\mu)\) 。

\(\nu\) 的唯一性由局部化原理（Ch. III, §2, No.~1, 命题 1 的推论）立即得出。为建立存在性，我们注意到：若 \(t \in \mathbf{T}\) ，则每个点 \(x \in \overline{\pi}^{-1}(t)\) 有一个邻域，它与 \(\overline{\pi}^{-1}(t)\) 只交于点 \(x\) ，从而 \(\overline{\pi}^{-1}(t)\) 是 \(\mathbf{X}\) 的一个离散子空间，且族 \((\varepsilon_x)_{x \in \overline{\pi}^{-1}(t)}\) 是可和的；我们用 \(\lambda_t\) 表示它的和。接着我们证明，映射 \(t \mapsto \lambda_t\) 是标量本质 \(\mu\) -可积的，且其积分 \(\nu = \int \lambda_t d\mu(t)\) 就是所求的逆像。这将由下列引理立即得出：

引理. --- a) 设 \(f\) 是 \(\mathcal{K}_{+}(X)\) 的一个元素；函数 \(t \mapsto \lambda_{t}(f)\) 是正的、上半连续的、具有紧支集，且它在 \(\pi(X)\) 上的限制是连续的。

\begin{enumerate}
\def\labelenumi{\alph{enumi})}
\setcounter{enumi}{1}
\tightlist
\item
  设 U 是 \(\mathcal{U}\) 的一个元素，\(\nu\) 是标量本质 \(\mu\) -可积函数 \(t \mapsto \lambda_{t}\) 的积分；测度 \(\nu|U\) 在 \(\pi|U\) 下的像等于 \(\mu|\pi(U)\) 。
\end{enumerate}

为建立 a)，可借助单位分解（Ch. III, §1, No.~2, 引理 1）归结到 f 的支集 S 含于一个开集 \(U \in \mathcal{U}\) 的情形。设 g 是映射 \(t \mapsto \lambda_{t}(f)\) ；由于 \(\pi|U\) 是同胚，\(g|\pi(U)\) 属于 \(\mathcal{K}_{+}(\pi(U))\) ，于是（\((\pi(U)\) 是 \(\pi(X)\) 中的开集）g 在 \(\pi(X)\) 上的限制是连续的。由于 g 是正的且 g 在紧集 \(\pi(S)\) 上的限制连续，可见 g 在 T 上是上半连续的。由此推出 g 是 \(\mu\) -可积的。

为建立 b)，用 g 表示 \(\mathcal{K}\big(\pi(\mathrm{U})\big)\) 的一个元素，用 \(g^{\circ}\) 表示它在 T 上的以 0 延拓，用 f 表示函数 \(g \circ (\pi|U)\) ，并用 \(f^{\circ}\) 表示 f 在 X 上的以 0 延拓。断言 b) 等价于等式 \(\int g^{\circ} d\mu = \int f^{\circ} d\nu\) 。但 \(f \in \mathcal{K}(\mathrm{U})\) ，因此 \(f^{\circ} \in \mathcal{K}(\mathrm{X})\) ，于是第二个积分等于 \(\int \lambda_{t}(f^{\circ}) d\mu(t)\) 。最后 \(\lambda_{t}(f^{\circ}) = g^{\circ}(t)\) ，这就完成了证明。

现在我们注意到 \(\pi(\mathrm{X})\) 在 T 中是开集，从而是 \(\mu\) -可测的；映射 \(\Lambda: t \mapsto \lambda_t\) 是淡 \(\mu\) -可测的，因为它在 \(\pi(\mathrm{X})\) 和 \(\mathbb{C}\pi(\mathrm{X})\) 中每一个上的限制都是淡连续的。在这些条件下，公式 \(\overline{\pi}^{1}(\mu) = \int \lambda_t d\mu(t)\) 定义了 \(\overline{\pi}^{1}(\mu)\) 关于 \(\pi\) 的一个切片分解，且命题 10 给出下列结果：

命题 11. --- 设 \(\pi\) 是局部紧空间 \(X\) 到局部紧空间 \(T\) 中的一个局部同胚，设 \(\mu\) 是 \(T\) 上的一个正测度。设 \(n\) 是这样一个数值函数：它对每个 \(t \in T\) 结合 \(\overline{\pi}^{1}(t)\) 的元素个数（若此数有限），而在相反情形结合 \(+\infty\) 。为使 \(\pi\) 是 \(\overline{\pi}^{1}(\mu)\) -逆紧的，必须且只须 \(n\) 是局部 \(\mu\) -可积的，此时

\[
\pi \big (\bar {\pi} ^ {1} (\mu) \big) = n \cdot \mu .\tag{15}
\]

